\section{Substitution and evaluation of power series}

\subsection{Defining substitution}

\begin{definition}
\label{def.fps.subs}
\lean{AlgebraicCombinatorics.fps_comp, AlgebraicCombinatorics.fps_comp_coeff, AlgebraicCombinatorics.fps_comp_coeff_finite}
\leantarget
\leanok
Let $f$ and $g$ be two FPSs in $K[[x]]$. Assume that $[x^0]g = 0$
(that is, $g = g_1 x + g_2 x^2 + g_3 x^3 + \cdots$ for some
$g_1, g_2, g_3, \ldots \in K$).

We then define an FPS $f[g] \in K[[x]]$ as follows:

Write $f$ in the form $f = \sum_{n \in \mathbb{N}} f_n x^n$ with
$f_0, f_1, f_2, \ldots \in K$. (That is, $f_n = [x^n]f$ for each
$n \in \mathbb{N}$.) Then, set
\[
f[g] := \sum_{n \in \mathbb{N}} f_n g^n.
\]
(This sum is well-defined, as we will see in Proposition
\ref{prop.fps.subs.wd} \textbf{(b)} below.)

This FPS $f[g]$ is also denoted by $f \circ g$, and is called the
\emph{composition} of $f$ with $g$, or the result of \emph{substituting}
$g$ for $x$ in $f$.

Equivalently, the $n$-th coefficient of $f[g]$ is the finite sum
\[
[x^n](f[g]) = \sum_{d=0}^{n} f_d \cdot [x^n](g^d).
\]
\end{definition}

\begin{lemma}
\label{lem.fps.subs.eq-subst}
\lean{AlgebraicCombinatorics.fps_comp_eq_subst}
\leanhelper
\leanok
The composition $f[g]$ defined above equals the standard
substitution operation on formal power series.
\end{lemma}

\begin{proof}
\leanok
By definition.
\end{proof}

\subsection{Well-definedness of substitution}

\begin{proposition}
\label{prop.fps.subs.wd}
\lean{AlgebraicCombinatorics.fps_subs_wd_firstCoeffs, AlgebraicCombinatorics.fps_subs_wd_summable, AlgebraicCombinatorics.fps_subs_wd_constCoeff}
\leantarget
\leanok
Let $f$ and $g$ be two FPSs in $K[[x]]$. Assume that $[x^0]g = 0$.
Write $f$ in the form $f = \sum_{n \in \mathbb{N}} f_n x^n$ with
$f_0, f_1, f_2, \ldots \in K$. Then:

\textbf{(a)} For each $n \in \mathbb{N}$, the first $n$ coefficients of the
FPS $g^n$ are $0$.

\textbf{(b)} The sum $\sum_{n \in \mathbb{N}} f_n g^n$ is well-defined, i.e.,
the family $(f_n g^n)_{n \in \mathbb{N}}$ is summable.

\textbf{(c)} We have $[x^0](\sum_{n \in \mathbb{N}} f_n g^n) = f_0$.
\end{proposition}

\begin{proof}
\textbf{(a)} The FPS $g$ has constant term $[x^0]g = 0$. Hence, by
Lemma~\ref{lem.fps.g=xh} (applied to $a = g$), there exists an
$h \in K[[x]]$ such that $g = xh$. Now, let $n \in \mathbb{N}$. From
$g = xh$, we obtain $g^n = (xh)^n = x^n h^n$. However,
Lemma~\ref{lem.fps.first-n-coeffs-of-xna} (applied to $k = n$ and
$a = h^n$) yields that the first $n$ coefficients of $x^n h^n$ are $0$.
In other words, the first $n$ coefficients of $g^n$ are $0$.

\textbf{(b)} This follows from part \textbf{(a)}. We must prove that for each
$n \in \mathbb{N}$, all but finitely many $i \in \mathbb{N}$ satisfy
$[x^n](f_i g^i) = 0$. Fix $n \in \mathbb{N}$. Let $i \in \mathbb{N}$
satisfy $i > n$. Then $n < i$, and by part \textbf{(a)} (applied to $i$
instead of $n$), the first $i$ coefficients of $g^i$ are $0$. Since $n < i$,
the coefficient $[x^n](g^i) = 0$. Thus $[x^n](f_i g^i) = f_i \cdot 0 = 0$.

\textbf{(c)} Let $n$ be a positive integer. By part \textbf{(a)}, the first
$n$ coefficients of $g^n$ are $0$, so $[x^0](g^n) = 0$ and thus
$[x^0](f_n g^n) = 0$. Therefore,
\[
[x^0]\left(\sum_{n \in \mathbb{N}} f_n g^n\right)
= [x^0](f_0 \underbrace{g^0}_{= 1}) + \sum_{n > 0} \underbrace{[x^0](f_n g^n)}_{= 0}
= f_0.
\]
\end{proof}

\begin{lemma}
\label{lem.fps.subs.zero}
\lean{AlgebraicCombinatorics.subst_zero}
\leanhelper
\leanok
Let $g \in K[[x]]$ with $[x^0]g = 0$. Then $0[g] = 0$.
\end{lemma}

\begin{proof}
\leanok
Immediate from the definitions: $0[g] = \sum_{n \in \mathbb{N}} 0 \cdot g^n = 0$.
\end{proof}

\begin{lemma}
\label{lem.fps.subs.one}
\lean{AlgebraicCombinatorics.subst_one}
\leanhelper
\leanok
Let $g \in K[[x]]$ with $[x^0]g = 0$. Then $\underline{1}[g] = \underline{1}$.
\end{lemma}

\begin{proof}
\leanok
Immediate from the definitions: $\underline{1}[g] = 1 \cdot g^0 + \sum_{n > 0} 0 \cdot g^n = 1$.
\end{proof}

\subsection{Lemma for first $k$ coefficients}

\begin{lemma}
\label{lem.fps.fg-coeffs-0}
\lean{AlgebraicCombinatorics.fps_fg_coeffs_zero}
\leantarget
\leanok
Let $f, g \in K[[x]]$ satisfy $[x^0]g = 0$. Let $k \in \mathbb{N}$ be such
that the first $k$ coefficients of $f$ are $0$. Then, the first $k$
coefficients of $f \circ g$ are $0$.
\end{lemma}

\begin{proof}
\leanok
We have $[x^0]g = 0$. Hence, by Lemma~\ref{lem.fps.g=xh} (applied to
$a = g$), there exists an $h \in K[[x]]$ such that $g = xh$.

Write $f = (f_0, f_1, f_2, \ldots)$. The first $k$ coefficients of $f$
are $0$, so $f_n = 0$ for each $n < k$. Now,
\[
f \circ g = \sum_{n \in \mathbb{N}} f_n g^n
= \sum_{\substack{n \in \mathbb{N}; \\ n < k}} \underbrace{f_n}_{= 0} g^n
+ \sum_{\substack{n \in \mathbb{N}; \\ n \geq k}} f_n \underbrace{g^n}_{= x^n h^n}
= \sum_{\substack{n \in \mathbb{N}; \\ n \geq k}} f_n x^n h^n.
\]
For any $m < k$, each term $f_n x^n h^n$ with $n \geq k$ satisfies
$[x^m](x^n h^n) = 0$ (since $m < k \leq n$). Thus $[x^m](f \circ g) = 0$
for each $m < k$.
\end{proof}

\subsection{Kronecker delta notation}

\begin{definition}
\label{def.kron-delta}
\lean{AlgebraicCombinatorics.kroneckerDelta}
\leantarget
\leanok
If $i$ and $j$ are any objects, then $\delta_{i,j}$ means the element
\[
\delta_{i,j} =
\begin{cases}
1, & \text{if } i = j; \\
0, & \text{if } i \neq j
\end{cases}
\]
of $K$.
\end{definition}

\begin{lemma}
\label{lem.kron-delta.self}
\lean{AlgebraicCombinatorics.kroneckerDelta_self}
\leanhelper
\leanok
$\delta_{i,i} = 1$.
\end{lemma}

\begin{proof}
\leanok
Immediate from the definition.
\end{proof}

\begin{lemma}
\label{lem.kron-delta.ne}
\lean{AlgebraicCombinatorics.kroneckerDelta_ne}
\leanhelper
\leanok
If $i \neq j$, then $\delta_{i,j} = 0$.
\end{lemma}

\begin{proof}
\leanok
Immediate from the definition.
\end{proof}

\begin{lemma}
\label{lem.kron-delta.comm}
\lean{AlgebraicCombinatorics.kroneckerDelta_comm}
\leanhelper
\leanok
$\delta_{i,j} = \delta_{j,i}$.
\end{lemma}

\begin{proof}
\leanok
Case split on whether $i = j$.
\end{proof}

\begin{lemma}
\label{lem.kron-delta.mul-left}
\lean{AlgebraicCombinatorics.kroneckerDelta_mul_left}
\leanhelper
\leanok
$\delta_{i,j} \cdot a = \begin{cases} a & \text{if } i = j, \\ 0 & \text{if } i \neq j. \end{cases}$
\end{lemma}

\begin{proof}
\leanok
Case split on whether $i = j$.
\end{proof}

\begin{lemma}
\label{lem.kron-delta.mul-right}
\lean{AlgebraicCombinatorics.kroneckerDelta_mul_right}
\leanhelper
\leanok
$a \cdot \delta_{i,j} = \begin{cases} a & \text{if } i = j, \\ 0 & \text{if } i \neq j. \end{cases}$
\end{lemma}

\begin{proof}
\leanok
Case split on whether $i = j$.
\end{proof}

\begin{lemma}
\label{lem.kron-delta.sum-left}
\lean{AlgebraicCombinatorics.sum_kroneckerDelta_left}
\leanhelper
\leanok
If $\alpha$ is a finite type, then $\sum_{i \in \alpha} \delta_{i,j} = 1$.
\end{lemma}

\begin{proof}
\leanok
All terms vanish except $i = j$, which contributes $1$.
\end{proof}

\begin{lemma}
\label{lem.kron-delta.sum-right}
\lean{AlgebraicCombinatorics.sum_kroneckerDelta_right}
\leanhelper
\leanok
If $\alpha$ is a finite type, then $\sum_{j \in \alpha} \delta_{i,j} = 1$.
\end{lemma}

\begin{proof}
\leanok
By symmetry ($\delta_{i,j} = \delta_{j,i}$), this reduces to
Lemma~\ref{lem.kron-delta.sum-left}.
\end{proof}

\begin{lemma}
\label{lem.kron-delta.contraction}
\lean{AlgebraicCombinatorics.sum_kroneckerDelta_mul}
\leanhelper
\leanok
If $\alpha$ is a finite type, then
$\sum_{i \in \alpha} \delta_{i,j} \cdot f(i) = f(j)$.
\end{lemma}

\begin{proof}
\leanok
By the multiplication property, $\delta_{i,j} \cdot f(i) = 0$ for $i \neq j$
and $f(j)$ for $i = j$. The sum collapses to the single term $f(j)$.
\end{proof}

\begin{lemma}
\label{lem.kron-delta.contraction-right}
\lean{AlgebraicCombinatorics.sum_mul_kroneckerDelta}
\leanhelper
\leanok
If $\alpha$ is a finite type, then
$\sum_{j \in \alpha} f(j) \cdot \delta_{i,j} = f(i)$.
\end{lemma}

\begin{proof}
\leanok
Analogous to Lemma~\ref{lem.kron-delta.contraction}.
\end{proof}

\begin{lemma}
\label{lem.kron-delta.nat-eq}
\lean{AlgebraicCombinatorics.kroneckerDelta_nat_eq}
\leanhelper
\leanok
For natural numbers $n,m\in\mathbb{N}$, we have
$\delta_{n,m} = \begin{cases} 1 & \text{if } n = m, \\ 0 & \text{if } n \neq m. \end{cases}$
\end{lemma}

\begin{proof}
\leanok
This is an immediate specialization of the definition to natural numbers.
\end{proof}

\begin{lemma}
\label{lem.kron-delta.int-eq}
\lean{AlgebraicCombinatorics.kroneckerDelta_int_eq}
\leanhelper
\leanok
For integers $n,m\in\mathbb{Z}$, we have
$\delta_{n,m} = \begin{cases} 1 & \text{if } n = m, \\ 0 & \text{if } n \neq m. \end{cases}$
\end{lemma}

\begin{proof}
\leanok
This is an immediate specialization of the definition to integers.
\end{proof}

\begin{lemma}
\label{lem.kron-delta.eq-zero-iff}
\lean{AlgebraicCombinatorics.kroneckerDelta_eq_zero_iff}
\leanhelper
\leanok
For a nontrivial ring $K$ and any $i,j$, we have
$\delta_{i,j} = 0$ if and only if $i \neq j$.
\end{lemma}

\begin{proof}
\leanok
If $i \neq j$, then $\delta_{i,j} = 0$ by definition.
Conversely, if $i = j$, then $\delta_{i,j} = 1 \neq 0$ (since $K$ is nontrivial).
\end{proof}

\begin{lemma}
\label{lem.kron-delta.eq-one-iff}
\lean{AlgebraicCombinatorics.kroneckerDelta_eq_one_iff}
\leanhelper
\leanok
For a nontrivial ring $K$ and any $i,j$, we have
$\delta_{i,j} = 1$ if and only if $i = j$.
\end{lemma}

\begin{proof}
\leanok
If $i = j$, then $\delta_{i,j} = 1$ by definition.
Conversely, if $i \neq j$, then $\delta_{i,j} = 0 \neq 1$ (since $K$ is nontrivial).
\end{proof}

\subsection{Laws of substitution}

\begin{proposition}
\label{prop.fps.subs.rules}
\lean{AlgebraicCombinatorics.fps_subs_add, AlgebraicCombinatorics.fps_subs_mul, AlgebraicCombinatorics.fps_subs_div, AlgebraicCombinatorics.fps_subs_pow, AlgebraicCombinatorics.fps_subs_assoc_constCoeff, AlgebraicCombinatorics.fps_subs_assoc, AlgebraicCombinatorics.fps_subs_const, AlgebraicCombinatorics.fps_subs_X_left, AlgebraicCombinatorics.fps_subs_X_right, AlgebraicCombinatorics.fps_subs_sum, AlgebraicCombinatorics.fps_subs_summable, AlgebraicCombinatorics.fps_subs_summableFPSSum}
\leantarget
\leanok
Composition of FPSs satisfies the rules you would expect it to satisfy:

\textbf{(a)} If $f_1, f_2, g \in K[[x]]$ satisfy $[x^0]g = 0$, then
$(f_1 + f_2) \circ g = f_1 \circ g + f_2 \circ g$.

\textbf{(b)} If $f_1, f_2, g \in K[[x]]$ satisfy $[x^0]g = 0$, then
$(f_1 \cdot f_2) \circ g = (f_1 \circ g) \cdot (f_2 \circ g)$.

\textbf{(c)} If $f_1, f_2, g \in K[[x]]$ satisfy $[x^0]g = 0$, then
$\dfrac{f_1}{f_2} \circ g = \dfrac{f_1 \circ g}{f_2 \circ g}$, as long as
$f_2$ is invertible. (In particular, $f_2 \circ g$ is automatically invertible
under these assumptions.)

\textbf{(d)} If $f, g \in K[[x]]$ satisfy $[x^0]g = 0$, then
$f^k \circ g = (f \circ g)^k$ for each $k \in \mathbb{N}$.

\textbf{(e)} If $f, g, h \in K[[x]]$ satisfy $[x^0]g = 0$ and $[x^0]h = 0$,
then $[x^0](g \circ h) = 0$ and $(f \circ g) \circ h = f \circ (g \circ h)$.

\textbf{(f)} We have $\underline{a} \circ g = \underline{a}$ for each $a \in K$
and $g \in K[[x]]$.

\textbf{(g)} We have $x \circ g = g \circ x = g$ for each $g \in K[[x]]$.

\textbf{(h)} If $(f_i)_{i \in I} \in K[[x]]^I$ is a summable family of FPSs,
and if $g \in K[[x]]$ is an FPS satisfying $[x^0]g = 0$, then the family
$(f_i \circ g)_{i \in I} \in K[[x]]^I$ is summable as well and we have
$(\sum_{i \in I} f_i) \circ g = \sum_{i \in I} f_i \circ g$.
\end{proposition}

\begin{proof}
\textbf{(a)} Write $f_1 = \sum_{n \in \mathbb{N}} f_{1,n} x^n$ and
$f_2 = \sum_{n \in \mathbb{N}} f_{2,n} x^n$. Then
$f_1 + f_2 = \sum_{n \in \mathbb{N}} (f_{1,n} + f_{2,n}) x^n$, so
\[
(f_1 + f_2)[g] = \sum_{n \in \mathbb{N}} (f_{1,n} + f_{2,n}) g^n
= \sum_{n \in \mathbb{N}} f_{1,n} g^n + \sum_{n \in \mathbb{N}} f_{2,n} g^n
= f_1[g] + f_2[g].
\]

\textbf{(f)} We have $\underline{a} = \sum_{n \in \mathbb{N}} a \delta_{n,0} x^n$.
Hence $\underline{a}[g] = \sum_{n \in \mathbb{N}} a \delta_{n,0} g^n
= a \cdot 1 + 0 = \underline{a}$.

\textbf{(g)} We have $x = \sum_{n \in \mathbb{N}} \delta_{n,1} x^n$, so
$x[g] = \sum_{n \in \mathbb{N}} \delta_{n,1} g^n = g^1 = g$. Also,
writing $g = \sum_{n \in \mathbb{N}} g_n x^n$, we get
$g[x] = \sum_{n \in \mathbb{N}} g_n x^n = g$.

\textbf{(b)} Write $f_1 = \sum_{i \in \mathbb{N}} f_{1,i} x^i$ and
$f_2 = \sum_{j \in \mathbb{N}} f_{2,j} x^j$. Then
\[
f_1[g] \cdot f_2[g]
= \left(\sum_{i \in \mathbb{N}} f_{1,i} g^i\right)
  \left(\sum_{j \in \mathbb{N}} f_{2,j} g^j\right)
= \sum_{n \in \mathbb{N}} \left(\sum_{\substack{(i,j) \in \mathbb{N}^2; \\ i+j=n}} f_{1,i} f_{2,j}\right) g^n.
\]
The same computation with $g$ replaced by $x$ gives
$f_1 \cdot f_2 = \sum_{n \in \mathbb{N}} c_n x^n$ where
$c_n = \sum_{\substack{(i,j) \in \mathbb{N}^2; \\ i+j=n}} f_{1,i} f_{2,j}$.
Thus $(f_1 \cdot f_2)[g] = \sum_{n \in \mathbb{N}} c_n g^n = f_1[g] \cdot f_2[g]$.

The interchange of infinite summation signs (discrete Fubini) is justified because
the family $(f_{1,i} f_{2,j} g^{i+j})_{(i,j) \in \mathbb{N}^2}$ is summable:
for any $m \in \mathbb{N}$ and any $(i,j)$ with $m < i+j$, we have
$[x^m](g^{i+j}) = 0$ by Proposition~\ref{prop.fps.subs.wd} \textbf{(a)}.

\textbf{(c)} Assume $f_2$ is invertible. By part \textbf{(b)} applied to $f_2^{-1}$,
$(f_2^{-1} \cdot f_2) \circ g = (f_2^{-1} \circ g) \cdot (f_2 \circ g)$, so
$(f_2^{-1} \circ g) \cdot (f_2 \circ g) = \underline{1} \circ g = \underline{1}$
by part \textbf{(f)}. Hence $f_2 \circ g$ is invertible. Then by part
\textbf{(b)} applied to $f_1/f_2$,
$(f_1/f_2 \cdot f_2) \circ g = (f_1/f_2 \circ g) \cdot (f_2 \circ g)$, i.e.,
$f_1 \circ g = (f_1/f_2 \circ g) \cdot (f_2 \circ g)$. Dividing by $f_2 \circ g$
gives $f_1/f_2 \circ g = (f_1 \circ g)/(f_2 \circ g)$.

\textbf{(d)} By induction on $k$. Base case: $f^0 \circ g = \underline{1} \circ g
= \underline{1} = (f \circ g)^0$ by part \textbf{(f)}. Induction step:
$f^{m+1} \circ g = (f \cdot f^m) \circ g = (f \circ g) \cdot (f^m \circ g)
= (f \circ g) \cdot (f \circ g)^m = (f \circ g)^{m+1}$ by part \textbf{(b)}.

\textbf{(h)} First, we show $(f_i \circ g)_{i \in I}$ is summable. Since
$(f_i)_{i \in I}$ is summable, for each $n \in \mathbb{N}$ there is a finite
$I_n \subseteq I$ such that $[x^n] f_i = 0$ for all $i \in I \setminus I_n$.
For $i \in I \setminus (I_0 \cup I_1 \cup \cdots \cup I_n)$, the first $n+1$
coefficients of $f_i$ are all $0$, so by Lemma~\ref{lem.fps.fg-coeffs-0},
$[x^n](f_i \circ g) = 0$. Thus the family $(f_i \circ g)_{i \in I}$ is summable.

The equality $(\sum_{i \in I} f_i) \circ g = \sum_{i \in I} f_i \circ g$ follows by
writing each $f_i = \sum_{n \in \mathbb{N}} f_{i,n} x^n$ and interchanging
summation signs (justified by summability of the family
$(f_{i,n} g^n)_{(i,n) \in I \times \mathbb{N}}$).

\textbf{(e)} Write $g = \sum_{n \in \mathbb{N}} g_n x^n$. By
Proposition~\ref{prop.fps.subs.wd} \textbf{(c)} (applied to $g, h$),
$[x^0](g \circ h) = g_0 = [x^0]g = 0$.

For associativity: $f \circ g = \sum_{n \in \mathbb{N}} f_n g^n$, and the
family $(f_n g^n)_{n \in \mathbb{N}}$ is summable by
Proposition~\ref{prop.fps.subs.wd} \textbf{(b)}. By part \textbf{(h)},
\[
(f \circ g) \circ h
= \left(\sum_{n \in \mathbb{N}} f_n g^n\right) \circ h
= \sum_{n \in \mathbb{N}} (f_n g^n) \circ h.
\]
By parts \textbf{(b)}, \textbf{(d)}, and \textbf{(f)},
$(f_n g^n) \circ h = (\underline{f_n} \cdot g^n) \circ h
= (\underline{f_n} \circ h) \cdot (g^n \circ h)
= \underline{f_n} \cdot (g \circ h)^n
= f_n (g \circ h)^n$.
Hence $(f \circ g) \circ h = \sum_{n \in \mathbb{N}} f_n (g \circ h)^n
= f \circ (g \circ h)$.
\end{proof}

\subsection{Example: Fibonacci generating function}

\begin{definition}
\label{def.fps.geometric-series}
\lean{AlgebraicCombinatorics.geometricSeries}
\leanhelper
\leanok
The \emph{geometric series} is the FPS $\frac{1}{1-x} = (1-x)^{-1} \in K[[x]]$
(where $K$ is a field).
\end{definition}

\begin{lemma}
\label{lem.fps.geometric-subst-fibonacci}
\lean{AlgebraicCombinatorics.fps_geometric_subst_fibonacci}
\leanhelper
\leanok
Substituting $x + x^2$ into the geometric series $\frac{1}{1-x}$ yields
the generating function for shifted Fibonacci numbers:
\[
\frac{1}{1-x}\Big[x + x^2\Big] = \frac{1}{1 - x - x^2}.
\]
\end{lemma}

\begin{proof}
By Proposition~\ref{prop.fps.subs.rules} \textbf{(c)}, substitution preserves
inverses when the original has invertible constant coefficient. Since
$[x^0](1-x) = 1 \neq 0$, we have
$\frac{1}{1-x}\big[x+x^2\big] = \frac{1}{(1-x)[x+x^2]}$. Now
$(1-x)[x+x^2] = 1 - (x+x^2) = 1 - x - x^2$ by the linearity of substitution,
so the result follows.
\end{proof}
